\cotefile{74}
\scanpage{89}
et les \inserted{deux} structures \uline{co} produit obtenues en mettant sur l'un des facteurs une structure de doublet, sur l'autre une structure de segment. Celle en \deleted{\unsure{évidence|induite}} dessinée est
\diagram{\begin{tikzcd}[ampersand replacement=\&,row sep=small,column sep=large]
(ba') \arrow[r,dash] \arrow[dr,dash] \& (bb') \arrow[dl,dash,crossing over]\\
(aa') \arrow[r,dash] \& (ab')
\end{tikzcd}}

On trouve ainsi \uline{les} trois \inserted{uniques} structures de carré sur l'ens à 4 él. produit, \unsure{correspondt|correspondant} aux partitions de cet \uline{ens} de type $(2,2)$\ldots.\par\centerline{\rule{1.1cm}{.4pt}}

Soit un carré $C$, défini par l'ens $C$ de ses sommets, et la partition de type $2,2$, $(C_\alpha)_{\alpha\in\underline D}$ de $C$ en les deux diagonales\inserted{(ou une involution \unsure{ss pt fixe|ss pts fixes|sans pt fixe})}. Une orientation du carré définit, pour chaque $\alpha\in\underline D$.
$\alpha=\{x,x'\}$, une bijection \inserted{$\varphi_\alpha:z\mapsto uz$} de $C_\alpha$ \deleted{$=$} $\{x,x'\}$ sur $C_\beta=\{y,y'\}$. On a $\varphi_\beta\varphi_\alpha=-\mathrm{id}_{C_\alpha}$, $\varphi_\alpha\varphi_\beta=-\mathrm{id}_{C_\beta}$
\diagram{\begin{tikzcd}[ampersand replacement=\&,row sep=small,column sep=small]
x \arrow[r,"\overset{u}{\longrightarrow}"] \arrow[d,dash] \arrow[dr,dash] \& y \arrow[d,dash] \arrow[dl,dash]\\
y' \arrow[r,dash] \& x'
\end{tikzcd}}
(chacune de ces relations impliquant l'autre). D'où
\[
\mathrm{or}(C)\simeq
(\underbrace{\mathop{\bigwedge}_{\alpha\in\underline D} C_\alpha}_{\substack{\underline D^{*}=\text{ens des diagonales}\\\text{du carré dual}}})\wedge\underline D
\simeq \underline D^{*}\wedge\underline D
\]
% archive sheet number: 88
